\chapter{Conclusion}

This project presented ANON4REID, a systematic benchmark for evaluating the privacy--utility trade-off in person re-identification on the Market-1501 dataset. Four anonymization techniques of increasing strength were applied across all 36,036 images: Gaussian blur with MediaPipe face detection (L1), Stable Diffusion inpainting (L2), Realistic Anonymization by Diffusion with ControlNet and LCM-LoRA (L3), and Canny edge detection (L4).

Utility was measured by the \acrshort{reid} performance of a pre-trained \acrshort{osnet} model, and privacy was measured by the identity classification accuracy of a MobileNetV2-based adversarial attacker trained on the original test images. The results reveal a stark trade-off:

\begin{keybox}[Summary of Results]
\begin{itemize}[noitemsep]
    \item Face-only blurring (L1) retains 80.9\% of baseline mAP but provides almost no privacy protection, with the attacker still achieving 88.66\% Top-1 accuracy, compared to a random-chance baseline of 0.13\%.
    \item Body-level anonymization (L2--L4) achieves strong privacy protection ($\geq$88.24\%) but reduces \acrshort{reid} utility to below 2.3\% of baseline.
    \item No intermediate ``sweet spot'' was found: the privacy--utility relationship exhibits a cliff rather than a smooth curve, with no evaluated technique achieving both meaningful privacy and usable \acrshort{reid} performance simultaneously.
\end{itemize}
\end{keybox}

These findings demonstrate that in the context of person \acrshort{reid} on low-resolution pedestrian imagery, the features that enable cross-camera matching and the features that enable identity recognition are fundamentally overlapping. Effective anonymization necessarily destroys re-identification capability, and preserving re-identification capability necessarily leaves identity information exposed to adversarial classifiers.

This has practical implications for the deployment of privacy-preserving surveillance systems under GDPR and similar regulatory frameworks. Face blurring alone is insufficient against learned adversaries that exploit body-level features, and stronger anonymization techniques that do provide robust privacy protection render \acrshort{reid} functionality inoperative. Organizations seeking to deploy \acrshort{reid} systems in privacy-sensitive contexts must acknowledge this fundamental tension rather than assuming that simple face anonymization provides adequate protection.

The complete evaluation pipeline, including anonymization scripts, attacker training, utility and privacy evaluation, and an interactive Streamlit dashboard, is provided as open-source software to support reproducibility and future research.
